\section{Proofs for stopping accessibility and certified iteration}\label{app:accessibility}
\begin{proof}[Proof of Proposition~\ref{prop:traceconsistency}]
Fix $x_0\in[1.5,2)$ and use $c=.8$, $\theta=p=0$. Wealth decreases, stays below two, and satisfies $x_s\geq x_0-.8s\geq.7$ for $s\leq1$. Thus it does not hit either wealth face. Preferences equal $2+.05W_s$. On the event $H=\{\sup_{s\leq1}|.05W_s|\leq.2\}$ there is no preference exit, $u\in[1.8,2.2]$, and running utility is at least $-1.5$. To see the last inequality, set $r=u-1\in[.8,1.2]$. The function $-(5/4)^r/r$ is increasing on this interval since $r\log(5/4)<1$. Its value at $.8$ exceeds $-1.5$, using $(5/4)^4<(6/5)^5$. Settlement on $H$ is at least $-.05$, because $-.02(.2)^2+.1\log(.7)>-.05$. Discounting negative terms cannot worsen these lower bounds. The pathwise discounted payoff on $H$ is therefore at least $-1.55$.

Until the original stopping time, $u\in[1.2,2.8]$. The same elementary monotonicity bounds running utility by $-6$ from below, and $G\geq-.1$. Even on $H^c$, the stopped payoff is at least $-6-8-.1=-14.1$. The Brownian maximal inequality gives $\Pr(H^c)\leq2e^{-8}$. Hence
\[
 J_2(0,2,x_0;.8,0,0)\geq-1.55-12.55(2e^{-8}).
\]
Subtracting $g(0,2,2)=.1\log2-8$ proves \eqref{eq:tracejump}, uniformly as $x_0\uparrow2$. Directed MPFR evaluation encloses its right-hand side in the interval recorded in the replication; its lower endpoint exceeds $6.37$.

Finally, a smooth witness equal to $g$ on all stopping faces and with $\sup_a\cR_av\leq e_+(s)$ gives $V(0,z)\leq v(0,z)+\int_0^1e^{-.04s}e_+(s)ds$ by the stopped It\^o identity. Take $z=(2,x)$, let $x\uparrow2$, and use continuity of $v$ and the established strict gap. The integrated uniform positive residual must exceed $6.37$. A sequence with vanishing such residuals is excluded, without making any assertion that the economic control problem lacks a solution.
\end{proof}

\begin{proof}[Proof of Theorem~\ref{thm:accessibility}]
Discounted It\^o's formula, localization of the stochastic integral, and the bounded stopped-domain hypotheses give, for every admissible control $a$,
\[
 J(t,z;a)-v(t,z)=\E_{t,z}^a\!\left[\int_t^\tau e^{-\rho(s-t)}\cR_av(s,Z_s)ds
 +e^{-\rho(\tau-t)}(g-v)(\tau,Z_\tau)\right].
\]
The positive residual bound and the facewise upper discrepancy bounds imply
\[
 V-v\leq\int_t^T e^{-\rho(s-t)}e_+(s)ds+\sum_i\beta_i^+\omega_i^*.
\]
Apply the same identity to $\pi$ with the opposite sign to obtain
\[
 v-J^\pi\leq\int_t^T e^{-\rho(s-t)}e_-(s)ds+\sum_i\beta_i^-\omega_i^\pi.
\]
Adding proves the assertion. No common stopping time across two distinct controlled paths is assumed. For the proposed smooth hitting-probability bound $w_i$, optional stopping yields $\E[e^{-\rho(\tau-t)}w_i(\tau,Z_\tau)]\leq w_i(t,z)$; domination of the indicator on the stopping boundary then gives the stated probability bound.
\end{proof}

\begin{proof}[Proof of Proposition~\ref{prop:architecture}]
Write $Y=2-x>0$ and $p=Yq$, where $q=(.15+.65\tanh N_p)/1.5$ is bounded. Until another exit,
\[
 dY=(c-.04)ds+Y\{.02-.06xq\}ds-.2xqY\,dZ^x_s.
\]
The bounded coefficients define a strictly positive stochastic exponential $\Phi$. Variation of constants gives
\[
 Y_s=\Phi_s\left(Y_t+\int_t^s\Phi_r^{-1}(c_r-.04)dr\right)>0,
\]
since $c\geq.05$. The upper wealth face is therefore inaccessible from the interior. On the remaining state faces $B=0$; at terminal time $1-t=0$. On the upper wealth face $B\geq0$ and softplus is positive, so $v-g\geq0$. Thus every upper-comparison trace discrepancy is zero, and the only potentially positive policy discrepancy has zero hitting probability. The same calculation applies to admissible bounded clipped head perturbations before the distance multiplication.
\end{proof}

\begin{proof}[Proof of Theorem~\ref{thm:iteration}]
Let $d_n=\epsilon_n/(1-q)$. The resolvent of $I-qP_{\pi_n}$ is nonnegative and has row sums at most $(1-q)^{-1}$, so $\|v_n-V^{\pi_n}\|_\infty\leq d_n$. Put $a_n=2qd_n+\eta_n$. The contraction inequalities and approximate greedy inequality give
\[
 TV^{\pi_n}-T_{\pi_{n+1}}V^{\pi_n}\leq a_n\mathbf1.
\]
Since $TV^{\pi_n}\geq V^{\pi_n}$, the same nonnegative resolvent gives $V^{\pi_{n+1}}\geq V^{\pi_n}-a_n\mathbf1/(1-q)$. Therefore
\begin{align*}
 V^*-V^{\pi_{n+1}}
 &=TV^*-T_{\pi_{n+1}}V^{\pi_{n+1}}\\
 &\leq qE_n\mathbf1+a_n\mathbf1+qP_{\pi_{n+1}}(V^{\pi_n}-V^{\pi_{n+1}})\\
 &\leq\{qE_n+a_n/(1-q)\}\mathbf1.
\end{align*}
Taking the supremum norm proves \eqref{eq:iterationrate}. The rate follows by summing the resulting geometric convolution. The comparison with a continuous-time problem is the triangle inequality after separately bounding the two transfer terms; it does not follow from the finite Markov contraction alone.
\end{proof}

\section{Proof of the state-global inventory certificate}\label{app:inventory}
\begin{proof}[Proof of Theorem~\ref{thm:inventory}]
The scalar backward Riccati solutions exist on the finite horizon, remain at least $.5$, and are bounded above. Indeed, in remaining time their drift is $2a_i p-p^2+q_i$, positive at $.5$ and negative above its positive equilibrium. The neural gains also remain at least $.5$ by construction. Let $d_i=\widehat p_i-p_i^*$ and let $R_i$ be the neural Riccati residual. Then
\[
 d_i'=(\widehat p_i+p_i^*-2a_i)d_i+R_i,\qquad d_i(1)=0.
\]
The coefficient is at least $\kappa_i=1-2a_i$. Integrating backward therefore yields $|d_i(t)|\leq\varepsilon_i(1-e^{-\kappa_i(1-t)})/\kappa_i\leq D$.

The quadratic optimal value has matrix $P^*=p_\perp^*(I-M)+p_m^*M$ and a scalar noise term. Completing the square in its generator and applying It\^o's formula gives the exact identity
\[
 \mathcal J(0,x;\widehat a)-\mathcal V(0,x)
 =\E\int_0^1|\{\widehat P(s)-P^*(s)\}X_s|^2ds.
\]
The two projectors are orthogonal, so the integrand is bounded by $D^2|X_s|^2$. Both eigenvalues of $A-\widehat P$ are at most $-.4$. It\^o's formula for $|X_s|^2$ then gives
\[
 \frac{d}{ds}\E|X_s|^2\leq-.8\E|X_s|^2+.01d.
\]
Gronwall's inequality and $|x|^2/d\leq1$ imply $d^{-1}\E|X_s|^2\leq e^{-.8s}+(.01/.8)(1-e^{-.8s})$. Integration proves \eqref{eq:inventorycertificate}. The bounded deterministic gains give finite second moments and justify localization and removal of the stopping sequence in both It\^o identities. No bounded-state truncation or simulated tail probability is used.
\end{proof}

\section{Proof of the compensation bound}
\begin{proof}[Proof of Proposition~\ref{prop:topup}]
Holding the specified paths fixed, the derivative of consumption utility with respect to the top-up is
\[
 \frac{\partial}{\partial\lambda}\frac{\{(1+\lambda)c\}^{1-u}}{1-u}
 =c^{1-u}(1+\lambda)^{-u}\geq(1+\lambda)^{-2.8}.
\]
The inequality uses $c\leq1$, $u>1$, $u\leq2.8$, and $\lambda\geq0$. Integrating from zero to $\lambda$, then over the discounted alive-time measure, bounds the welfare gain below by $A\{1-(1+\lambda)^{-1.8}\}/1.8$. Solving for a gain of $\epsilon$ proves \eqref{eq:topup}. For the time-control library, wealth feasibility excludes a wealth exit, and the retained drift bounds and Brownian maximal inequality give the displayed preference-exit probability. Surviving to terminal time alone yields the stated lower bound on discounted alive time. This argument requires no extrapolation of the certified optimal value to an altered wealth state.
\end{proof}
